\subsection{Same-frequency model rejection versus positive-$D$ detection}

A small deterministic fit residual relative to the full channel vector is not by itself a statistical model-acceptance statement. We therefore place the same-frequency one-mode goodness-of-fit test and the positive-$D$ test on the same explicit theoretical noise scale. At one RF frequency the six complex channels provide 12 real observations, while the complex $(C,K,r)$ one-mode model has six real fitted parameters, leaving six local residual degrees of freedom.

We retain the independent equal real/imaginary quadrature-noise model used below and define $S=\sqrt{\langle |J_m|^2\rangle}/\sigma_{\rm q}$. For the one-mode test, the deterministic post-fit channel residual supplies the noncentral alternative. For positive-$D$ detectability, the full $12\times6$ channel Jacobian is propagated through the fitted complex root and then through $D(\gamma)$; an idealized one-sided Gaussian $D>0$ test is evaluated at the same $\alpha=0.0027$ and 90\% power.

\begin{table}[t]
\caption{Same-frequency statistical ordering under the reference independent-quadrature noise model. $S_D$ is the RMS-channel SNR required for 90\% power to establish positive apparent diffusion; $S_{1\mathrm{m}}$ is the SNR required for 90\% power to reject the six-channel one-mode manifold, both at $\alpha=0.0027$.}
\label{tab:samefreq}
\centering
\begin{tabular}{rrrr}
\toprule
RF & $D_{\rm eff}$ (m$^2$/s) & $S_D$ (dB) & $S_{1\mathrm{m}}$ (dB)\\
\midrule
100 MHz & 2.610e-03 & 111.86 & 102.29\\
500 MHz & 2.549e-03 & 70.04 & 88.19\\
1000 MHz & 2.349e-03 & 52.41 & 81.80\\
\bottomrule
\end{tabular}
\end{table}

Under this reference covariance, the ordering is frequency dependent.  At the tested point(s) 100 MHz: $S_{1\mathrm{m}}=102.29$ dB before $S_D=111.86$ dB, so the spectral-model check self-announces before positive apparent diffusion reaches the stated 90\% detection power.  At the remaining tested point(s), positive apparent diffusion reaches 90\% power first, leaving finite RMS-channel-SNR windows (500 MHz: $70.04$--$88.19$ dB; 1 GHz: $52.41$--$81.80$ dB) in which $D>0$ has reached the stated detection power while rejection of the same-frequency one-mode manifold has not.  The example therefore supports conditional same-frequency hidden-risk windows at those RF points, not a universal stealth claim.

This comparison also clarifies the role of the deterministic full-vector residual quoted elsewhere in the paper: it is a descriptive approximation error, whereas Table~\ref{tab:samefreq} is the corresponding covariance-aware model-selection statement for the reference noise model.
